\section{Conclusion}
\label{sec:conclusion}

Honeypots have spent two decades recording attacker behavior without
reasoning about it, and the arrival of generative models promised to close
that gap while opening a worse one: a deception tool whose product is
unverified improvisation. This paper argued that the right unit of design
is not the generator but the \emph{loop} around it. \system couples
behavioral profiling, retrieval-grounded planning, a fiction graph, and
serving-time verification into a cycle in which every response is either
verified against the session's deception state or served from vetted
deterministic content, with a human gate on the highest-stakes material
and an audit trail over every decision.

The evaluation supports the division of labor the architecture claims,
and the escalation analysis prices it rung by rung. Retrieval grounding
is the rung the system cannot lose: removing
it costs $+0.154$ composite ($[+0.092, +0.221]$) and every detected TTP,
while leaving persona realism untouched, a null result the design predicts.
The fiction graph and verifier maintain their properties under targeted
adversarial pressure: zero contradictions across 45 re-queries, zero
unverified serves across a 218-query live campaign, and an ablation
battery in which each mechanism's removal costs the property it was added
to provide. The deception holds a three-evaluator consensus of $3.86/5$ with
$0.750$ exact agreement across personas of one judge model -- a
within-model consistency measure rather than inter-annotator agreement
-- and two independent full campaigns reproduce the
composite to $+0.005$ with the deterministic strata replicating exactly.
At the top of the ladder, the result is a subtraction:
kind-awareness fixed 12/13 real-log assessments and lifted bait exposure
by 125\%, added no measurable dwell time, and showed no marginal routing
effect that the benchmark could resolve -- so it left the serving path.
Each rung we retained paid for itself in its measured contribution, and
the rung we removed is the result that carries furthest beyond this
system.

The limits are stated with the same directness. The deployment hears no ambient
internet; the judge shares a provider with the generator; and the
availability of generative components is structurally hostage to provider
quota and alignment posture, a constraint the fallback path absorbs but
does not remove. None of these limits is hidden by the loop; all of them
are measured by it. We release the scenario sets, the scoring harnesses,
and every persisted artifact behind the numbers in this paper -- code,
cloud stack, results, and the compiled paper itself -- at
\url{https://github.com/0xzodiac/The-Intelligence-Ladder}, so that the
claims can be replayed rather than believed.

We see three directions as the most productive next steps. First, extending the fiction
graph across sessions, so that the longitudinal attacker memory and the
consistency machinery stop being separate systems and cross-session
deception inherits cross-session verification. Second, field evaluation on
internet-reachable deployments, where the ambient-traffic axis this
testbed excludes by scope choice becomes measurable, and the steering
hypothesis, better bait, measurable attacker dwell, can be tested against
adversaries who did not consent to be studied. Third, treating the ladder
audit as an ongoing protocol: as the deterministic machinery matures,
each intelligent rung should be re-priced against its cheapest
certifiable competitor, and the loop should shrink accordingly. The
system this paper releases is one where that audit has already run once,
and the serving path is smaller because of it.
